% =============================================================================
% 空 agent 集成用紧凑模板：南京理工大学学士论文 LaTeX 格式
% 覆盖所有结构组件，最小化正文内容，专注格式复用
% 编译: tectonic main.tex  或  xelatex && bibtex && xelatex && xelatex
% =============================================================================
\documentclass[a4paper, 12pt, twoside]{ctexart}

% --- 页面设置 (A4; 上30mm、下24mm、左25mm、右25mm; 页眉footskip=20mm) ---
\usepackage[
  a4paper,
  top=30mm, bottom=24mm, left=25mm, right=25mm,
  headheight=15pt, headsep=10mm, footskip=20mm
]{geometry}

% --- 正文：小四号宋体，固定约20磅行距 ---
\usepackage{setspace}
\AtBeginDocument{\zihao{-4}\songti\setlength{\baselineskip}{20pt}\setlength{\parindent}{2em}}

% --- 防止 Overfull \hbox（允许段落适当拉伸词间距以避免溢出） ---
\emergencystretch=0.5em

% --- 字体设置 ---
% 中文：宋体为主体；标题处使用宋体伪粗以匹配学校"宋体加粗"要求
% 英文与数字：Times New Roman
\usepackage{fontspec}
\setmainfont{Times New Roman}
\setCJKmainfont[AutoFakeBold=2.5]{Songti SC}
\setCJKfamilyfont{song}[AutoFakeBold=2.5]{Songti SC}
\setCJKfamilyfont{hei}{Heiti SC}

% --- 数学公式 ---
\usepackage{amsmath, amssymb}

% --- 图表 ---
\usepackage{graphicx}
\graphicspath{{../图表/}}
\usepackage{float}
\usepackage{subcaption}
\DeclareCaptionFont{songwu}{\songti\zihao{5}}
\captionsetup{font=songwu, labelfont=songwu, labelsep=space}
\captionsetup[lstlisting]{font=songwu, labelfont=songwu, labelsep=space}
\usepackage{booktabs}
\usepackage{multirow}
\usepackage{longtable}
\usepackage{tabularx}
\usepackage{adjustbox}
\usepackage{afterpage}

% --- 图表按章编号（图X-X、表X-X、式X-X）---
\IfFileExists{chngcntr.sty}{\usepackage{chngcntr}}{}
\counterwithin{figure}{section}
\counterwithin{table}{section}
\counterwithin{equation}{section}
\def\thefigure{\thesection.\arabic{figure}}
\def\thetable{\thesection.\arabic{table}}
\def\theequation{\thesection.\arabic{equation}}

% --- 代码块 ---
\usepackage{listings}
\usepackage{listingsutf8}
\usepackage{xcolor}
\lstset{
    basicstyle=\zihao{5}\ttfamily,
    columns=fullflexible,
    keepspaces=true,
    literate={，}{{，}}1 {。}{{。}}1 {：}{{：}}1 {；}{{；}}1
            {（}{{（}}1 {）}{{）}}1 {、}{{、}}1,
    keywordstyle=\color{blue},
    commentstyle=\color{gray},
    stringstyle=\color{red!70!black},
    numbers=left,
    numberstyle=\tiny\color{gray},
    frame=single,
    breaklines=true,
    tabsize=4,
}

% --- 交叉引用与超链接（hyperref 应尽量靠后加载）---
\usepackage{hyperref}
\hypersetup{
    colorlinks=true,
    linkcolor=black,
    citecolor=black,
    urlcolor=black,
}

% --- 参考文献（GB/T 7714-2015 数字编号，sort&compress 压缩编号）---
% 使用 natbib + gbt7714-numerical；如用 BibTeX 则xelatex两遍+bibtex+两遍
\usepackage[numbers,sort&compress]{natbib}
\bibliographystyle{gbt7714-numerical}

% --- 页眉页脚（奇偶页不同，页码外侧）---
% 偶数页(CE)：左→章节标题，右→"学士学位论文"
% 奇数页(CO)：左→"学士学位论文"，右→章节标题
% 页码位于 LE（左偶）、RO（右奇）即外侧
\usepackage{fancyhdr}
\pagestyle{fancy}
\fancyhf{}
% 正确的 sectionmark 定义（必须放在使用之前）
\makeatletter
\def\sectionmark#1{\markboth{\thesection\quad #1}{}}
\makeatother
\fancyhead[CO]{\songti\zihao{5} 学士学位论文\hfill <论文题目(简)>}
\fancyhead[CE]{\songti\zihao{5} \leftmark\hfill 学士学位论文}
\fancyfoot[LE,RO]{\songti\zihao{5}\thepage}
\fancypagestyle{plain}{%
  \fancyhf{}%
  \fancyhead[CO]{\songti\zihao{5} 学士学位论文\hfill <论文题目(简)>}%
  \fancyhead[CE]{\songti\zihao{5} \leftmark\hfill 学士学位论文}%
  \fancyfoot[LE,RO]{\songti\zihao{5}\thepage}%
}

% --- 目录深度（含一二三级标题）---
% tocloft: \cftsecnumdepth=3 即含 section/sub/subsub 三级
\IfFileExists{tocloft.sty}{\usepackage[titles]{tocloft}}{}
\def\contentsname{目\quad 录}
\def\refname{参考文献}

% --- 章节标题格式（学校规范）---
\ctexset{
    section = {
        format = \raggedright\songti\bfseries\zihao{-3},
        beforeskip = 18pt,
        afterskip = 18pt,
        break = \clearpage,
    },
    subsection = {
        format = \raggedright\songti\bfseries\zihao{4},
        beforeskip = 12pt,
        afterskip = 12pt,
    },
    subsubsection = {
        format = \raggedright\songti\bfseries\zihao{-4},
        beforeskip = 6pt,
        afterskip = 6pt,
    },
    paragraph = {
        format = \raggedright\songti\zihao{-4},
        beforeskip = 6pt,
        afterskip = 0pt,
    },
}

% --- 自定义命令 ---
\newcommand{\frontmattertitle}[1]{%
  \vspace*{18pt}\begin{center}{\songti\bfseries\zihao{3}#1}\end{center}\vspace{18pt}}
\newcommand{\centeredsectionstar}[1]{%
  \clearpage\phantomsection\vspace*{18pt}\begin{center}{\songti\bfseries\zihao{3}#1}\end{center}\vspace{18pt}}

% =============================================================================
\begin{document}

% ========================== 摘要 ==========================
\newpage
\pagestyle{plain}
\pagenumbering{roman}

\frontmattertitle{摘\quad 要}
\phantomsection
\addcontentsline{toc}{section}{摘要}

本文针对XXXXXXXXXXXXXXXXXXXXXXXXXXXXXXXXXXXXXXXXXXXXXXXXXXXXXXXXXXXXXXXXXXXXXXXXXX。
XXXXXXXXXXXXXXXXXXXXXXXXXXXXXXXXXXXXXXXXXXXXXXXXXXXXXXXXXXXXXXXXXXXXXXXXXXXXXXXXXXX。
XXXXXXXXXXXXXXXXXXXXXXXXXXXXXXXXXXXXXXXXXXXXXXXXXXXXXXXXXXXXXXXXXXXXXXXXXXXXXXXXXXX。

\vspace{1cm}
\noindent{\songti\bfseries\zihao{4} 关键词：}{\songti\zihao{-4} 关键词1；关键词2；关键词3}

% ========================== Abstract ==========================
\newpage
\markboth{Abstract}{}
\begin{center}{\bfseries\zihao{3} Abstract}\end{center}\vspace{18pt}
\phantomsection
\addcontentsline{toc}{section}{Abstract}

This thesis presents a systematic study of XXXXXXXXXXXXXXXXXXXXXXXXXXXXXXXXXXXXXXXX.
XXXXXXXXXXXXXXXXXXXXXXXXXXXXXXXXXXXXXXXXXXXXXXXXXXXXXXXXXXXXXXXXXXXXXXXXXXXXXXXXXXX.
XXXXXXXXXXXXXXXXXXXXXXXXXXXXXXXXXXXXXXXXXXXXXXXXXXXXXXXXXXXXXXXXXXXXXXXXXXXXXXXXXXX.

\vspace{1cm}
\noindent{\bfseries\zihao{4} Keywords:}{\zihao{-4} keyword1; keyword2; keyword3}

% ========================== 目录 ==========================
\newpage
\markboth{目\quad 录}{}
\begingroup
\def\contentsname{目\quad 录}
\tableofcontents
\endgroup
% 双面打印时，若前置部分总页数为奇数则补空白页
\ifodd\value{page}\newpage\thispagestyle{empty}\mbox{}\fi

% ========================== 正文 ==========================
\newpage
\pagestyle{fancy}
\pagenumbering{arabic}

% ========================================================================
%   第一章  绪论
% ========================================================================
\section{绪论}

\subsection{研究背景与意义}

XXXXXXXXXXXXXXXXXXXXXXXXXXXXXXXXXXXXXXXXXXXXXXXXXXXXXXXXXXXXXXXXXXXXXXXXXXXXX。
如图\ref{fig:architecture}所示，XXXXXXXXXXXXXXXXXXXXXXXXXXXXXXXXXXXXXXXXX。

\begin{figure}[H]
    \centering
    \includegraphics[width=0.65\textwidth]{architecture.jpg}
    \caption{系统架构图}
    \label{fig:architecture}
\end{figure}

\subsection{国内外研究现状}

XXXXXXXXXXXXXXXXXXXXXXXXXXXXXXXXXXXXXXXXXXXXXXXXXXXXXXXXXXXXXXXXXXXXXXXXXXXXX。
文献\cite{author2024example}提出了XXXXXXXXXXXXXXXXXXXXXXXXXXXXXXXXXXXXXXXXXX。

\subsubsection{Linux调度器研究}

XXXXXXXXXXXXXXXXXXXXXXXXXXXXXXXXXXXXXXXXXXXXXXXXXXXXXXXXXXXXXXXXXXXXXXXXXXXXX。

% 第四级标题使用 paragraph（学校规范通常不含 subsubsubsection）
\paragraph{相关技术基础}

XXXXXXXXXXXXXXXXXXXXXXXXXXXXXXXXXXXXXXXXXXXXXXXXXXXXXXXXXXXXXXXXXXXXXXXXXXXX。

\subsection{本文创新点}

本文主要创新点如下：

\begin{enumerate}
    \item 创新点一：XXXXXXXXXXXXXXXXXXXXXXXXXXXXXXXXXXXXXXXXXXXXXXXX。
    \item 创新点二：XXXXXXXXXXXXXXXXXXXXXXXXXXXXXXXXXXXXXXXXXXXXXXXX。
\end{enumerate}

% ========================================================================
%   第二章  相关技术基础
% ========================================================================
\section{相关技术基础}

本章介绍XXXXXXXXXXXXXXXXXXXXXXXXXXXXXXXXXXXXXXXXXXXXXXXXXXXXXXXXXXXXXXXX。

\subsection{Linux调度器原理}

CFS的核心思想是XXXXXXXXXXXXXXXXXXXXXXXXXXXXXXXXXXXXXXXXXXXXXXXXXXXXXXXXX。
虚拟运行时间的计算如式\eqref{eq:vruntime}所示。

\begin{equation}
    \Delta vruntime = \Delta t_{wall} \times \frac{NICE\_0\_LOAD}{weight(task)}
    \label{eq:vruntime}
\end{equation}

表\ref{tab:cfs-params}列出了CFS的核心可调参数及其语义。

\begin{table}[H]
    \centering
    \caption{CFS核心可调参数}
    \label{tab:cfs-params}
    \begin{tabular}{lll}
        \toprule
        \textbf{参数名称} & \textbf{基准值} & \textbf{语义} \\
        \midrule
        \texttt{sched\_latency\_ns} & 6\,ms & 调度周期目标 \\
        \texttt{sched\_min\_granularity\_ns} & 0.75\,ms & 最小抢占粒度 \\
        \bottomrule
    \end{tabular}
\end{table}

% ========================================================================
%   第三章  系统设计
% ========================================================================
\section{系统设计}

XXXXXXXXXXXXXXXXXXXXXXXXXXXXXXXXXXXXXXXXXXXXXXXXXXXXXXXXXXXXXXXXXXXXXXXXXXXXX。

\subsection{总体架构}

XXXXXXXXXXXXXXXXXXXXXXXXXXXXXXXXXXXXXXXXXXXXXXXXXXXXXXXXXXXXXXXXXXXXXXXXXXXXX。

\section{实验与分析}

XXXXXXXXXXXXXXXXXXXXXXXXXXXXXXXXXXXXXXXXXXXXXXXXXXXXXXXXXXXXXXXXXXXXXXXXXXXXX。

% ========================================================================
%   参考文献
% ========================================================================
% 注意：参考文献页码跟随正文继续使用阿拉伯数字，不切换回罗马数字
% 如使用 BibTeX：将 \begin{thebibliography}...替换为
%   \bibliography{refs}  并运行 xelatex && bibtex && xelatex && xelatex
\newpage
\begin{thebibliography}{99}
  \bibitem{author2024example} Author A, Author B. Title of the reference[J].
    Journal Name, 2024, 1(1): 1-10. DOI: 10.0000/example.
  \bibitem{linuxkernel} Linux Kernel Documentation[EB/OL].
    \url{https://www.kernel.org/doc/html/latest/}.
\end{thebibliography}

% ========================== 附录 ==========================
\appendix
\centeredsectionstar{附\quad 录}
\markboth{附\quad 录}{}
\addcontentsline{toc}{section}{附录}

\subsection*{附录A：实验配置}

表~\ref{tab:appendix-configs}列出了实验配置。

\begin{table}[H]
    \centering
    \caption{实验配置表}
    \label{tab:appendix-configs}
    \begin{tabular}{llr}
        \toprule
        \textbf{参数} & \textbf{符号} & \textbf{值} \\
        \midrule
        CPU核心数 & $N$ & 4 \\
        调度周期 & $T$ & 10\,ms \\
        \bottomrule
    \end{tabular}
\end{table}

\subsection*{附录B：核心代码}

\begin{lstlisting}[language=C, caption={示例代码}, label={lst:sample}]
#include <stdio.h>
int main() {
    printf("Hello, World!\n");
    return 0;
}
\end{lstlisting}

\end{document}
